% ----------------------------------------------------------
% Introdução
% ----------------------------------------------------------
\chapter{Introdução}

ERP’s, sistemas de auxílio à tomada de decisão ou websites de venda digital são exemplos de sistemas do meio digital desenvolvidos para melhorar aspectos de uma empresa ou possibilitar canais antes não utilizados. Todo tipo de sistema de informação é fruto de um projeto que foi planejado, desenvolvido, testado e implementado, e este também é um projeto que segue a mesma natureza.

\section{Caracterização do Projeto}
O presente documento evidencia e engloba etapas, diagramas e características deste projeto de software, que é alinhado também aos parâmetros de projeto acadêmico definidos pela norma ABNT \cite{de2011associaccao}, tendo como finalidade documentar todo o processo de desenvolvimento do sistema baseado no contexto de um negócio externo à instituição bem como registrar o andamento do Projeto Integrador de Extensão do curso superior de graduação de ADS.

\section{Objetivos}
Os objetivos do projeto indicam as ambições finais das atividades desenvolvidas e são divididos em um objetivo geral e alguns objetivo específicos.

\subsection{Objetivo Geral}
O objetivo principal deste trabalho é desenvolver um software que ajude a atender uma necessidade real de uma organização terceira ou comunidade além de envolver sua gestão interna.

\subsection{Objetivos Específicos}
Os objetivos específicos que foram definidos abrangem desde a necessidade do parceiro até à forma de desenvolvimento.
\begin{itemize}
  \item Levantar e documentar os requisitos funcionais e não funcionais;
  \item Disponibilizar relatórios que auxiliem o acompanhamento das operações e do desempenho do negócio;
  \item Desenvolver funcionalidades para cadastro, consulta, atualização e gerenciamento do acervo de livros;
  \item Implementar mecanismos para controle e acompanhamento do estoque;
  \item Implementar recursos de recomendação de livros com base no histórico de compras dos clientes;
\end{itemize}

\section{Justificativa}
O parceiro realiza atividades relacionadas à comercialização de livros usados, demandando um controle eficiente sobre seu acervo, estoque e operações de venda. A ausência de uma solução integrada pode gerar dificuldades na organização das informações, na consulta de títulos disponíveis e no acompanhamento das movimentações do negócio. Dessa forma, a criação de um sistema de gestão torna-se relevante por permitir a centralização dos dados, a redução de processos manuais, a melhoria na confiabilidade das informações e o aumento da eficiência operacional do estabelecimento. Além dos benefícios para o parceiro, o projeto proporciona aos estudantes a oportunidade de aplicar na prática os conhecimentos adquiridos ao longo do curso, promovendo a integração entre ensino, tecnologia e extensão universitária.

\section{Escopo do Projeto}
As atividades desenvolvidas incluem o desenvolvimento de toda a plataforma web, suas funcionalidades(registro de estoque, de clientes, histórico), documentação e andamento do projeto. No entanto, são excluídas do escopo a implementação prática e toda a infraestrutura necessária para tal.

\section{Estrutura do Documento}
O presente documento segue a fórmula de projeto acadêmico manuscrito determinado pela ABNT \cite{de2011associaccao} com diversas modificações para se alinhar aos propósitos e visão do curso de ADS. Para tanto, sua estrutura está dividida em basicamente 6 partes principais: a introdução, o contexto do projeto, análise do sistema, projeto do sistema, implementação e conclusão. Respectivamente, o projeto é introduzido e então serão discutidos a forma que como foi elaborado, estruturado, detalhado, recursos que serão necessários para sua implementação e, por fim, observações finais do mesmo.